\nonstopmode{}
\documentclass[a4paper]{book}
\usepackage[times,hyper]{Rd}
\usepackage{makeidx}
\makeatletter\@ifl@t@r\fmtversion{2018/04/01}{}{\usepackage[utf8]{inputenc}}\makeatother
% \usepackage{graphicx} % @USE GRAPHICX@
\makeindex{}
\begin{document}
\chapter*{}
\begin{center}
{\textbf{\huge Package `leafblower'}}
\par\bigskip{\large \today}
\end{center}
\ifthenelse{\boolean{Rd@use@hyper}}{\hypersetup{pdftitle = {leafblower: High-Performance Survey Calibration via iEPPA and L-BFGS-B}}}{}
\ifthenelse{\boolean{Rd@use@hyper}}{\hypersetup{pdfauthor = {Dennis Alexis Valin Dittrich}}}{}
\begin{description}
\raggedright{}
\item[Title]\AsIs{High-Performance Survey Calibration via iEPPA and L-BFGS-B}
\item[Version]\AsIs{0.1.0}
\item[Author]\AsIs{Dennis Alexis Valin Dittrich [aut, cre]}
\item[Maintainer]\AsIs{Dennis Alexis Valin Dittrich }\email{davdittrich@gmail.com}\AsIs{}
\item[Description]\AsIs{Drop-in replacement for autumn::harvest() implementing iEPPA
(Sinkhorn block-coordinate descent) and L-BFGS-B over the Deville-Sarndal
logit dual. Adds min_weight lower bound; shares a single C++17 core with
the Python bindings.}
\item[License]\AsIs{MIT + file LICENSE}
\item[Encoding]\AsIs{UTF-8}
\item[NeedsCompilation]\AsIs{yes}
\item[SystemRequirements]\AsIs{C++17 compiler, OpenMP (optional, for SIMD exp
vectorisation)}
\item[Suggests]\AsIs{autumn (>= 0.2.0), testthat (>= 3.0.0), bench, lhs,
DiceKriging, ggplot2, rprojroot, survey}
\item[RoxygenNote]\AsIs{7.3.3}
\end{description}
\Rdcontents{Contents}
\HeaderA{anesrake}{anesrake compatibility wrapper}{anesrake}
%
\begin{Description}
Maps anesrake parameter names to harvest(). Unknown params produce warnings.
\end{Description}
%
\begin{Usage}
\begin{verbatim}
anesrake(
  inputter,
  targets,
  weightvec = NULL,
  caseid = NULL,
  pctlim = 0.05,
  cap = 5,
  choosemethod = "rake",
  type = "pctlim",
  nlim = 500L,
  iterate = TRUE,
  threads = 1L,
  ...
)
\end{verbatim}
\end{Usage}
%
\begin{Arguments}
\begin{ldescription}
\item[\code{inputter}] Data frame (maps to \code{data}).

\item[\code{targets}] Named list of target proportions.

\item[\code{weightvec}] Starting weights or NULL.

\item[\code{caseid}] Ignored (row label not used by calibration).

\item[\code{pctlim}] Deprecated; maps to \code{convergence["pct"]}.

\item[\code{cap}] Upper weight cap; maps to \code{max\_weight}.

\item[\code{choosemethod}] "rake" or "nrake"; maps to \code{method="lbfgsb"} with warning.

\item[\code{type}] Ignored.

\item[\code{nlim}] Max iterations; maps to \code{max\_iterations}.

\item[\code{iterate}] Ignored (always iterates).

\item[\code{threads}] Ignored (single-threaded v1).

\item[\code{...}] Unknown params produce a warning then are ignored.
\end{ldescription}
\end{Arguments}
%
\begin{Value}
Same as \code{harvest()}.
\end{Value}
\HeaderA{design\_effect}{Kish design effect (1-argument) or Henry-Valliant (4-argument)}{design.Rul.effect}
%
\begin{Description}
When called with only \code{weights}, computes the Kish (1992) design effect:
\code{n * sum(w\textasciicircum{}2) / sum(w)\textasciicircum{}2}.
\end{Description}
%
\begin{Usage}
\begin{verbatim}
design_effect(weights, outcome = NULL, data = NULL, target = NULL)
\end{verbatim}
\end{Usage}
%
\begin{Arguments}
\begin{ldescription}
\item[\code{weights}] Numeric vector of calibrated weights.

\item[\code{outcome}] Numeric outcome vector (optional; 4-arg form only).

\item[\code{data}] Data frame used in calibration (optional; 4-arg form only).

\item[\code{target}] Named list of target proportions (optional; 4-arg form only).
\end{ldescription}
\end{Arguments}
%
\begin{Details}
When called with all four arguments, computes the Henry \& Valliant (2015)
calibration design effect: \code{var\_weighted(y) / var\_unweighted(y)}.
\end{Details}
%
\begin{Value}
Numeric scalar design effect.
\end{Value}
\HeaderA{diagnose\_weights}{Diagnose calibration quality}{diagnose.Rul.weights}
%
\begin{Description}
Returns a data frame comparing original and weighted marginals to targets.
\end{Description}
%
\begin{Usage}
\begin{verbatim}
diagnose_weights(data, target, weights)
\end{verbatim}
\end{Usage}
%
\begin{Arguments}
\begin{ldescription}
\item[\code{data}] Data frame used in calibration.

\item[\code{target}] Named list of target proportions (same format as harvest()).

\item[\code{weights}] Numeric vector of calibrated weights, length nrow(data).
\end{ldescription}
\end{Arguments}
%
\begin{Value}
Data frame with columns: variable, level, prop\_original, prop\_weighted,
target, error\_original, error\_weighted.
\end{Value}
\HeaderA{effective\_sample\_size}{Effective sample size}{effective.Rul.sample.Rul.size}
%
\begin{Description}
Effective sample size
\end{Description}
%
\begin{Usage}
\begin{verbatim}
effective_sample_size(weights)
\end{verbatim}
\end{Usage}
%
\begin{Arguments}
\begin{ldescription}
\item[\code{weights}] Numeric vector of calibrated weights.
\end{ldescription}
\end{Arguments}
%
\begin{Value}
Numeric scalar: n / design\_effect(weights).
\end{Value}
\HeaderA{get\_current\_miss}{Current calibration miss (matching autumn's export name)}{get.Rul.current.Rul.miss}
%
\begin{Description}
Current calibration miss (matching autumn's export name)
\end{Description}
%
\begin{Usage}
\begin{verbatim}
get_current_miss(data, target, weights)
\end{verbatim}
\end{Usage}
%
\begin{Arguments}
\begin{ldescription}
\item[\code{data}] Data frame.

\item[\code{target}] Named list of target proportions.

\item[\code{weights}] Numeric weight vector.
\end{ldescription}
\end{Arguments}
%
\begin{Value}
Named numeric vector of max absolute errors per variable.
\end{Value}
\HeaderA{harvest}{Generate calibrated weights (drop-in for autumn::harvest)}{harvest}
%
\begin{Description}
Generate calibrated weights (drop-in for autumn::harvest)
\end{Description}
%
\begin{Usage}
\begin{verbatim}
harvest(
  data,
  target,
  min_weight = 0,
  max_weight = 5,
  method = "ieppa",
  verbose = 0,
  max_iterations = 500,
  start_weights = NULL,
  attach_weights = TRUE,
  weight_column = "weights",
  convergence = list(),
  sor = list(auto = TRUE, omega_min = 0.3),
  bounds_mode = "cell",
  homotopy_levels = 1L,
  homotopy_start_factor = 1,
  homotopy_end_factor = 1,
  homotopy_budget_p = 0.5,
  scheduler = c("round_robin", "greedy"),
  eta_schedule = c("fixed", "tang_dynamic"),
  eta_start = 1,
  eta_end = 1,
  eta_schedule_power = 0.5,
  select_params = NULL,
  select_function = NULL,
  error_function = NULL,
  adaptive_order = NULL,
  enforce_mean = TRUE,
  accelerate = FALSE,
  add_na_proportion = FALSE,
  auto_collapse = FALSE,
  collapse_vars = NULL,
  target_map = NULL,
  design_weights = NULL,
  ...
)
\end{verbatim}
\end{Usage}
%
\begin{Arguments}
\begin{ldescription}
\item[\code{data}] A data frame containing columns to calibrate.

\item[\code{target}] A named list of named numeric vectors (variable -> proportions).

\item[\code{min\_weight}] Lower bound on weights. Default 0 (no lower bound).

\item[\code{max\_weight}] Upper bound on weights. Default 5.

\item[\code{method}] Calibration method. One of \code{"auto"} (default: iEPPA or
raking based on M\_cell/n ratio), \code{"ieppa"} (paper-faithful iEPPA),
\code{"raking"} (IPF + Dykstra box projection), \code{"lbfgsb"}
(L-BFGS-B on concave dual), \code{"sinkhorn"} (KL Bregman Dykstra),
\code{"greg"} (Newton QP, Deville-Sarndal 1992), \code{"chebyshev"}
(L-infinity LP via IPM), \code{"grake"} (normalized Chebyshev via IPM).

\item[\code{verbose}] Integer verbosity: 0=silent, 1=progress, 2=debug.

\item[\code{max\_iterations}] Maximum inner BCD iterations per outer step. Default 500.

\item[\code{start\_weights}] Starting weights vector or NULL (uniform).

\item[\code{attach\_weights}] If TRUE, return data frame with weights column. Default TRUE.

\item[\code{weight\_column}] Name of weight column. Default "weights".

\item[\code{convergence}] Named list controlling the stopping criterion. Accepted keys:
\begin{itemize}

\item{} \code{metric}: quality metric to monitor. One of \code{"max\_err"}
(default), \code{"mean\_err"}, \code{"kl"}, \code{"chi2"},
\code{"grake\_norm"}, \code{"l1\_weight"}.
\item{} \code{rule}: stopping rule. One of \code{"improvement"} (default) —
stop when metric improves by less than \code{tol} relative; \code{"threshold"} —
stop when metric falls below \code{tol} absolutely; \code{"plateau"} —
stop when metric changes less than \code{tol} over a window.
\item{} \code{tol}: tolerance value (default \code{0.001}).
\item{} \code{stop\_when}: \code{"any"} (default) or \code{"all"}.

\end{itemize}

Shorthand keys (each sets metric+rule+tol simultaneously):
\begin{itemize}

\item{} \code{pct}: sets \code{metric="l1\_weight"}, \code{rule="plateau"},
\code{tol=pct}. Example: \code{list(pct = 0.001)}.
\item{} \code{absolute}: sets \code{metric="max\_err"}, \code{rule="threshold"},
\code{tol=absolute}. Example: \code{list(absolute = 1e-4)}.
\item{} \code{improvement}: sets \code{metric="max\_err"}, \code{rule="improvement"},
\code{tol=improvement}. Example: \code{list(improvement = 0.001)}.

\end{itemize}

Shorthands and explicit keys may be combined; explicit keys override shorthand
defaults. \code{convergence = list()} uses the default (max\_err + improvement +
tol = 0.001). For \code{method="ieppa"}, the default metric is \code{kl}
(Sinkhorn-type algorithm minimizing KL divergence). For other methods: \code{max\_err}.

\strong{chi2 cross-solver note:} chi2 is not directly comparable
across methods. iEPPA uses unnormalized cell mass as \code{W\_total};
raking and lbfgsb use \code{n}. Use chi2 as a convergence criterion
within one method; do not compare values across methods.

\item[\code{sor}] Named list for SOR adaptive under-relaxation (iEPPA only).
\code{NULL} disables SOR. Keys:
\begin{itemize}

\item{} \code{auto}: logical, default \code{TRUE}.
\item{} \code{omega\_min}: lower bound on relaxation factor, default \code{0.3}.
\item{} \code{omega}: fixed relaxation factor (disables auto-adapt).
\item{} \code{burnin}: iterations before adaptation starts, default \code{20}.

\end{itemize}


\item[\code{bounds\_mode}] One of "cell" (default) or "unit". Controls per-observation
vs cell-aggregate bound enforcement.

\item[\code{homotopy\_levels}] Number of homotopy levels (default 1 = disabled). Values > 1
progressively tighten max\_weight from homotopy\_start\_factor to homotopy\_end\_factor
across n levels.

\item[\code{homotopy\_start\_factor}] Starting max\_weight multiplier (default 1.0 = no change).

\item[\code{homotopy\_end\_factor}] Ending max\_weight multiplier (default 1.0 = no change).

\item[\code{homotopy\_budget\_p}] Budget split fraction across homotopy levels (default 0.5).

\item[\code{scheduler}] Margin sweep scheduler: "round\_robin" (default) or "greedy"
(Greenkhorn priority).

\item[\code{eta\_schedule}] ALM penalty schedule: "fixed" (default) or "tang\_dynamic".

\item[\code{eta\_start}] Starting ALM penalty multiplier (default 1.0).

\item[\code{eta\_end}] Ending ALM penalty multiplier (default 1.0).

\item[\code{eta\_schedule\_power}] Power for Tang-eta schedule interpolation (default 0.5).

\item[\code{select\_params}] Ignored with verbose >= 2 note.

\item[\code{select\_function}] Ignored.

\item[\code{error\_function}] Ignored.

\item[\code{adaptive\_order}] Ignored.

\item[\code{enforce\_mean}] Ignored (retained for compatibility).

\item[\code{accelerate}] Ignored.

\item[\code{add\_na\_proportion}] Not supported in v1; raises error if TRUE.

\item[\code{auto\_collapse}] Not supported in v1; raises error if TRUE.

\item[\code{collapse\_vars}] Not supported in v1; raises error if TRUE.

\item[\code{target\_map}] Passed through for data-frame target format handling.

\item[\code{...}] Additional arguments ignored.
\end{ldescription}
\end{Arguments}
%
\begin{Value}
data frame with weights column if \code{attach\_weights=TRUE}, else a numeric
vector of length \code{n}. In both cases the object carries attributes:
\begin{description}

\item[\code{result}] Named list of solver diagnostics. Key fields:
\begin{itemize}

\item{} \code{status}: 0=converged, 1=max\_iter hit, 2=infeasible, 3=bad args.
\item{} \code{iterations}: number of outer iterations completed.
\item{} \code{max\_error}: maximum marginal error at the returned weights.
\item{} \code{l1\_weight\_change}: L1-normalized weight change Σ|Δw|/n.
For \code{method="lbfgsb"} this is start-to-final (batch solver);
for iEPPA and raking it is iteration-to-iteration.
\item{} \code{grake\_norm}: survey-grake normalized residual
max\_k|misfit|/(1+|pop|).
\item{} \code{convergence\_used}: Named list with \code{metric}, \code{rule},
\code{tol}, and \code{fired\_at\_iter} documenting which criterion fired.
\item{} \code{best\_error}: minimum marginal error seen across all iterates.
\item{} \code{best\_weights}: numeric vector (length \code{n}, sum normalized to
\code{n}) at the iterate with minimum observed marginal error. When
\code{best\_error} is \code{Inf} (solver exited before the first convergence
check), \code{best\_weights} is all-zero. Guard:
\code{if (is.finite(attr(r, "result")\$best\_error))} before use.

\end{itemize}


\item[\code{algorithm}] Character name of the solver used.
\item[\code{iterations}] Convenience alias for \code{result\$iterations}.

\end{description}

\end{Value}
\HeaderA{weighted\_pct}{Weighted proportions}{weighted.Rul.pct}
%
\begin{Description}
Weighted proportions
\end{Description}
%
\begin{Usage}
\begin{verbatim}
weighted_pct(x, weights)
\end{verbatim}
\end{Usage}
%
\begin{Arguments}
\begin{ldescription}
\item[\code{x}] Factor or character vector.

\item[\code{weights}] Numeric weights, same length as x.
\end{ldescription}
\end{Arguments}
%
\begin{Value}
Named numeric vector of weighted proportions summing to 1.
\end{Value}
\printindex{}
\end{document}
